% SMCEP v1.0: standalone LaTeX methodological note with explicit in-text citations.
% Exploratory synthetic evidence; not a validation of real-world recursive self-improvement.
\documentclass[11pt,a4paper]{article}
\usepackage[utf8]{inputenc}
\usepackage[T1]{fontenc}
\usepackage{lmodern}
\usepackage{amsmath,amssymb,booktabs,geometry,hyperref}
\geometry{margin=2.6cm}
\title{SMCEP v1.0: Multi-Generation Policy Modification in a Synthetic Agent Environment}
\author{Research methods note}
\date{November 2026}
\begin{document}
\maketitle
\section{Motivation and relationship to prior work}
Self-modifying agents have been investigated through formal proof-oriented designs such as the G\"odel Machine \cite{schmidhuber2006}, as well as empirical code-modifying systems such as SICA \cite{robeyns2025}, the Darwin G\"odel Machine \cite{zhang2026}, and recursively improving AI research agents \cite{srikanth2026}. These works motivate a narrower measurement question: does a particular change to an agent's development policy improve its subsequent ability to produce independently validated improvements? The present experiment does not establish a general answer; it evaluates a specified synthetic mechanism.

SMCEP is studied as an independent account of self-modification and development-capacity transfer. Its governance implications intersect with, but are not subsumed by, the external authorization and verification mechanisms proposed in The Mind Genesis Architecture \cite{donmez2026}. This study does not evaluate that governance architecture.

\section{Design}
Two arms start from the same synthetic agent parameters. In the \emph{frozen-policy} arm, task parameters can be modified but proposal-policy parameters remain fixed. In the \emph{evolving-policy} arm, both can be modified. Both arms use the same candidate budget, task-generation rules, evaluator noise level, and matched random seeds. The synthetic task landscape includes nonlinear interactions and independent development, in-distribution (ID), and out-of-distribution (OOD) evaluation families. All held-out tasks are generated by the same researcher-defined landscape family; hence the transfer test is limited.

Let $X_t$ denote the current agent, $P(X_t)$ its task score, and $C_{\mathrm{ID}}(X_t;B)$ and $C_{\mathrm{OOD}}(X_t;B)$ the expected improvement produced by its development policy when optimizing a fixed reference agent under budget $B$. The policy-probe tasks are held out from the main selection tasks. The primary between-arm contrast is
\begin{equation}
\widehat\tau_Y=\frac{1}{N}\sum_{s=1}^{N}\left[(Y^{\mathrm{evolving}}_{T,s}-Y^{\mathrm{evolving}}_{0,s})-(Y^{\mathrm{frozen}}_{T,s}-Y^{\mathrm{frozen}}_{0,s})\right],
\end{equation}
where $Y\in\{P,C_{\mathrm{ID}},C_{\mathrm{OOD}}\}$ and $s$ indexes paired random seeds. This is a controlled contrast within the simulator, not an identified causal effect in real AI systems.

\section{Exploratory pilot results}
The pilot used $N=12$ matched seeds, $T=25$ generations, six candidates per generation, and twelve policy-probe repetitions. Values below are mean end-minus-start changes.
\begin{center}
\begin{tabular}{lrrr}
\toprule
Arm & $\Delta P$ & $\Delta C_{\mathrm{ID}}$ & $\Delta C_{\mathrm{OOD}}$\\
\midrule
Frozen policy & $+0.0730$ & $0.0000$ & $0.0000$\\
Evolving policy & $+0.0666$ & $-0.0071$ & $-0.0278$\\
\bottomrule
\end{tabular}
\end{center}
The paired evolving-minus-frozen differences were $-0.0064$ for task performance (approximate 95\% interval $[-0.0209,0.0080]$), $-0.0071$ for ID capacity ($[-0.0124,-0.0017]$), and $-0.0278$ for OOD capacity ($[-0.0498,-0.0057]$). Intervals are descriptive normal approximations across seeds, without multiplicity adjustment or preregistration.

\section{Interpretation and limitations}
The evolving-policy arm did not exhibit positive measured development-capacity transfer in this pilot. Zero capacity change in the frozen-policy arm is partly structural: the anchor-based probe measures policy parameters that are deliberately held constant in that arm. The evolving arm's lower OOD capacity therefore should not be interpreted as proof that recursive self-improvement is impossible. In addition, a single fixed synthetic landscape, limited sample size, and hand-specified proposal and selection mechanisms prevent generalization to real autonomous research agents. Follow-up work should include diverse task families, alternative proposal operators, independent holdouts, and pre-specified statistical comparisons.

\begin{thebibliography}{9}
\bibitem{schmidhuber2006} Schmidhuber, J. (2006). G\"odel Machines: Self-Referential Universal Problem Solvers Making Provably Optimal Self-Improvements. arXiv:cs/0309048v5. \url{https://arxiv.org/abs/cs/0309048}.
\bibitem{robeyns2025} Robeyns, M., Szummer, M., \& Aitchison, L. (2025). A Self-Improving Coding Agent. arXiv:2504.15228. \url{https://arxiv.org/abs/2504.15228}.
\bibitem{zhang2026} Zhang, J., Hu, S., Lu, C., Lange, R., \& Clune, J. (2026). Darwin G\"odel Machine: Open-Ended Evolution of Self-Improving Agents. ICLR 2026. arXiv:2505.22954. \url{https://arxiv.org/abs/2505.22954}.
\bibitem{srikanth2026} Srikanth, D., Zhao, B., Xu, D., Wu, Y., \& Jiang, Z. (2026). Recursive Self-Improvement of AI Research Agents. arXiv:2609.26457. \url{https://arxiv.org/abs/2609.26457}.
\bibitem{donmez2026} D\"onmez, A. (2026). The Mind Genesis Architecture: Shared Intelligence Governance at the Threshold of the ASI Age. SSRN preprint. \url{https://ssrn.com/abstract=7633039}.
\end{thebibliography}
\end{document}
